\section{Blanqueo del gradiente (Gradient Whitening)}

\subsection{Contexto del problema}

Al optimizar redes neuronales con SGD o AdamW, la observación empírica es que la mayoría de las matrices de parámetros 2D (pesos de las capas ocultas) tienen un número de condición alto, es decir, la matriz es de rango bajo. Esto hace que la actualización del gradiente tenga una magnitud muy grande en unas pocas direcciones, mientras que en una gran cantidad de «direcciones raras» la señal es sumamente débil.

\subsection{La operación central del blanqueo}

\begin{importantbox}{El núcleo de Muon: el blanqueo del gradiente}
Se aplica un blanqueo a la matriz del gradiente, \textbf{eliminando la magnitud/escala del gradiente y conservando solo la dirección}. Esto logra que:
\begin{itemize}
  \item se atenúe la influencia excesiva de las direcciones dominantes
  \item se amplifique la señal débil de las direcciones raras
  \item todas las direcciones obtengan una escala de actualización equitativa
\end{itemize}
\end{importantbox}

La operación de blanqueo se realiza mediante la SVD (descomposición en valores singulares):

$$
G = U \Sigma V^T \quad \xrightarrow{\text{blanqueo}} \quad \hat{G} = U V^T
$$

Es decir, se elimina la matriz de valores singulares $\Sigma$ y se conserva únicamente la parte de rotación $U$ y $V^T$.

\subsection{El significado geométrico de la SVD}

Para la matriz de gradiente $G$:
\begin{itemize}
  \item $U$: los vectores singulares izquierdos, que representan las direcciones del espacio de salida
  \item $\Sigma$: los valores singulares, que representan el «grado de importancia» (escala) de cada dirección
  \item $V^T$: los vectores singulares derechos, que representan las direcciones del espacio de entrada
\end{itemize}

Al eliminar $\Sigma$, $\hat{G} = UV^T$ resulta ser una \textbf{matriz ortogonal}, lo que garantiza que la magnitud de la actualización sea uniforme en todas las direcciones.

\subsection{Resumen de la sección}

El blanqueo del gradiente es la operación central de Muon: mediante la SVD se eliminan los valores singulares y se ortogonaliza la matriz de gradiente. Esto permite amplificar las señales raras que antes quedaban dominadas por las direcciones principales, lo que acelera el aprendizaje del modelo.
